\section{Encoding comparisons by locally reconstructible graphs}
\label{sec:encoding}

We now encode scalar comparisons between vector values at binary
addresses. A full-address value will be represented by a sum of $r$
entries, each depending on two consecutive address coordinates. The graph
will encode the table entries through local constraints, without listing
the exponentially many values of their sums.

The construction must also control graphs that were not produced by the
encoding. We will prove that every graph equivalent to the homogeneous
graph has the same constraint pieces, with possibly shifted right-hand
sides. An isomorphism will require one assignment to the shared table
entries satisfying all those shifts simultaneously.

\subsection{Comparisons and subdivided wires}

Throughout the construction we work over the field $\F$, fix
$r\ge 32$, and put $K=r+1$. Indices $i\in\{1,\ldots,r\}$ are cyclic, so
$r+1$ in an address-coordinate index means $1$. A finite nonempty
list of \emph{real types} is given. Each real type $P$ has dimension
$d_P\in\{1,2\}$; a value at type $P$ and address
$a\in\{0,1\}^r$ belongs to $\F^{d_P}$.

A \emph{comparison schema} $e$ consists of ordered real endpoints
$P,R$, a nonempty product domain
\[
D_e=\prod_{i=1}^r D_{e,i},
\qquad \varnothing\ne D_{e,i}\subseteq\{0,1\},
\]
coordinate injections $\sigma_{e,i}:D_{e,i}\to\{0,1\}$, and
scalar linear forms
\[
\lambda_{e,0}:\F^{d_P}\to\F,
\qquad
\lambda_{e,1}:\F^{d_R}\to\F.
\]
Write $\sigma_e(a)=(\sigma_{e,i}(a_i))_{i=1}^r$. The schema
represents comparisons between the first form at $(P,a)$ and the
second form at $(R,\sigma_e(a))$ for every $a\in D_e$. Zero
forms, repeated schemas, parallel schemas, and $P=R$ are all
allowed.

Replace each schema $e:P\to R$ by an oriented \emph{wire}
from $P$ to $R$ with $4K$ private internal types, each of dimension
one. Thus a wire has $4K+1$ links, and different wires have no
internal types in common. The real and internal types, joined by
these links, form the \emph{type layout}. For a type $t$, define its
allowed labels at coordinate $i$ by
\[
A_{t,i}=
\begin{cases}
\{0,1\},&t\text{ is real},\\
D_{e,i},&t\text{ is internal to wire }e.
\end{cases}
\]
Its dimension is denoted by $d_t$, extending the notation for real
types.

Write a link as $\ell=tu$, with $t$ first and $u$ second in its
wire's orientation. At coordinate $i$, its allowed label matches
are
\[
M_{\ell,i}=
\begin{cases}
\{(p,p):p\in D_{e,i}\},
 &\ell\text{ is not the last link of }e,\\
\{(p,\sigma_{e,i}(p)):p\in D_{e,i}\},
 &\ell\text{ is the last link of }e.
\end{cases}
\]
Give each endpoint of $\ell$ a scalar form $L_{\ell,t}$ or
$L_{\ell,u}$. At an internal endpoint this is the identity of
$\F$. At the first real endpoint of the wire it is
$\lambda_{e,0}$, and at the last real endpoint it is
$\lambda_{e,1}$. These assignments refer to the respective wire
ends even when both real endpoints are the same type.

\subsection{Readouts and blocks}

Introduce the following names for vector and scalar variables:
\begin{align*}
x_{i,t,p,p'}&\in\F^{d_t}
 &&(p\in A_{t,i},\ p'\in A_{t,i+1}),\\
y_{i,\ell,p,q}&\in\F
 &&((p,q)\in M_{\ell,i}).
\end{align*}
We call these variables \emph{readouts}; each vector readout has its
fixed ordered bit coordinates. It is useful to regard them as edges
between \emph{corners}, where a corner is a triple $(i,t,p)$ with
$p\in A_{t,i}$. The horizontal readout $x_{i,t,p,p'}$ joins
$(i,t,p)$ to $(i+1,t,p')$. The vertical readout
$y_{i,\ell,p,q}$ for $\ell=tu$ joins $(i,t,p)$ to $(i,u,q)$.

A full allowed address $a\in\prod_i A_{t,i}$ selects one horizontal
readout in each row. These edges form a cycle through the corners
$(i,t,a_i)$. Once values have been assigned to the horizontal names,
their sum
\[
                 h_t(a)=\sum_{i=1}^r x_{i,t,a_i,a_{i+1}}
\]
is the value represented at this address. It is determined by the
short tables $x_{i,t,\cdot,\cdot}$, not assigned independently at each
address. For example, the homogeneous comparison $e:P\to R$ is intended
to impose
\[
 \lambda_{e,0}(h_P(a))
     =\lambda_{e,1}(h_R(\sigma_e(a)))
       \qquad(a\in D_e).
\]
We implement this equality by local equations along the subdivided wire.
Their shifted versions will control the corresponding comparison in an
equivalent graph.

There are two kinds of \emph{blocks}. Each horizontal name $x$ has
one \emph{site block}, using just $x$. For each link $\ell=tu$,
row $i$, and matches
$(p,q)\in M_{\ell,i}$ and
$(p',q')\in M_{\ell,i+1}$, there is a \emph{square block} $B$
using the four names
\[
x_t=x_{i,t,p,p'},\qquad x_u=x_{i,u,q,q'},\qquad
y_i=y_{i,\ell,p,q},\qquad
y_{i+1}=y_{i+1,\ell,p',q'}.
\]
Figure~\ref{fig:square} records these shared names. A valuation of
a block assigns a value to each of its names, in that name's vector
or scalar space.

\begin{figure}[htbp]
\centering
\begin{tikzpicture}[x=1cm,y=1cm,every node/.style={font=\small}]
  \coordinate (tl) at (0,2);
  \coordinate (tr) at (5,2);
  \coordinate (bl) at (0,0);
  \coordinate (br) at (5,0);
  \draw (tl) -- node[midway,above] {$x_t$} (tr);
  \draw (bl) -- node[midway,below] {$x_u$} (br);
  \draw (tl) -- node[midway,left] {$y_i$} (bl);
  \draw (tr) -- node[midway,right] {$y_{i+1}$} (br);
  \fill (tl) circle (1.4pt);
  \fill (tr) circle (1.4pt);
  \fill (bl) circle (1.4pt);
  \fill (br) circle (1.4pt);
  \node[above left] at (tl) {$(i,t,p)$};
  \node[above right] at (tr) {$(i+1,t,p')$};
  \node[below left] at (bl) {$(i,u,q)$};
  \node[below right] at (br) {$(i+1,u,q')$};
\end{tikzpicture}
\caption{A square block at link $\ell=tu$ and row $i$.
Horizontal edges carry vector readouts at their types; vertical
edges carry scalar readouts at the matched labels. Its legal
valuations satisfy Equation~\eqref{eq:square}.}
\label{fig:square}
\end{figure}

Choose one bit $\beta_B\in\F$ for each square block. Every
valuation of a site is legal. A valuation of the displayed square
is legal if
\begin{equation}
\label{eq:square}
L_{\ell,t}(x_t)+L_{\ell,u}(x_u)+y_i+y_{i+1}=\beta_B.
\end{equation}
At least one endpoint of every link is internal, so a square uses
$d_t+d_u+2\le 5$ bit coordinates. The two vertical names are
distinct, and both have coefficient one. Consequently this is a
nonzero linear equation with at most sixteen legal valuations,
the same number for either right-hand side. This remains true when
an endpoint form is zero.

The long wires have a separate role in the game analysis. A collection
of at most $K$ blocks uses at most $K$ layout links, because a site uses
none and a square uses just one. It therefore cannot span a whole wire
of $4K+1$ links. Section~\ref{sec:projection} will use this separation
to extend assignments on small block collections.

The full address space has not been enumerated: each name or block uses
only one coordinate or two consecutive coordinates. We next turn these
bounded-size blocks into a graph in which any equivalent mate must retain
their incidence pattern.

\subsection{The uncolored graph}

First build a graph with indicated vertex colors.
\begin{enumerate}[label=\arabic*.]
\item For each bit coordinate of each readout name, make two
\emph{bit-value vertices}, labeled $0$ and $1$. Give this pair a
fresh color. The same pair is used in every block containing that
coordinate.
\item For each block, make one \emph{valuation vertex} for every
legal valuation, all with a fresh color specific to that block.
Connect a valuation vertex to the bit-value vertex specifying its
assigned value in each coordinate used by the block. Add no other
edges.
\end{enumerate}
This graph is bipartite, with valuation vertices on one side and
bit-value vertices on the other. Call these vertices the
\emph{base vertices}, and let their number be $N$. The number is
independent of the chosen bits $\beta_B$, by the preceding count.
Number the fresh colors, in one common order for all choices of
$\beta$, by $1,\ldots,J$.

To erase the colors, attach $j(N+1)$ new private leaves to each base
vertex of color $j$, and then forget the colors. The resulting
uncolored graph is denoted by $G(\beta)$. Here $\beta$ denotes the
list of all square bits, and $G(0)$ is the graph in which all these
bits vanish. The graph is simple and nonempty, and its order is
independent of $\beta$.

\begin{proposition}[Reconstruction of an arbitrary mate]
\label{prop:mate}
For any comparison data just described, with $r\ge 32$, every graph
$H$ satisfying $H\equiv_r G(0)$ is isomorphic to $G(\beta)$ for
some choice of one bit $\beta_B$ per square block.
\end{proposition}

\begin{proof}
The two graphs have the same positive order, and
Lemma~\ref{lem:wl-game} supplies a winning strategy from $G(0)$ to
$H$ with $K=r+1$ pair slots. We first recover the forgotten colors
and the private leaves, and then recover each block.

A base vertex of color $j$ in $G(0)$ has degree in the interval
\[
I_j=[j(N+1),\ j(N+1)+N-1].
\]
These intervals are pairwise disjoint and lie above degree one.
Every other vertex has degree one. By Lemma~\ref{lem:degree}, the
first announced bijection preserves degrees for every possible
vertex choice. It follows that $H$ has precisely the same number
of vertices in each $I_j$, the same number of degree-one vertices,
and no vertices of other degrees. Call its vertices in $I_j$ base
vertices of color $j$.

Any chosen base vertex of $H$ can be reached in the first
placement, by taking its inverse under the announced bijection.
The next announced bijection must preserve both degrees and
adjacency to that vertex. Its number of degree-one neighbors is
therefore exactly $j(N+1)$. Summing these required counts over
all base vertices accounts for the full number of degree-one
vertices in $H$. Since such a vertex has only one neighbor, all
are private leaves on bases; none remain in a separate component.
The degree intervals continue to preserve base colors throughout
every winning play.

There can be no base edge between two color classes whose vertices
are never adjacent in $G(0)$. Indeed, two arbitrary target vertices
can be reached through two successive placements, using the
inverse of each announced bijection. Their preimages have the
prescribed colors, and partial isomorphism preserves the edge
test. This also rules out edges within either side of the original
bipartition and edges to unrelated bit coordinates.

Fix a block. Mark every one of its valuation vertices in $G(0)$
and both bit-value vertices of each coordinate it uses. For a
square there are at most $16+2\cdot 5=26$ such vertices. A site
has at most four valuations and four bit-value vertices, so is
smaller. All these placements fit within $K\ge 33$ slots. Their
images are distinct, preserve their colors, and exhaust the
corresponding classes in $H$, since the class sizes agree. Thus
the entire colored incidence piece for this block is isomorphic
to its intended piece.

Label the two vertices in every bit pair of $H$ by $0$ and $1$,
once for the whole graph. In a local colored isomorphism, each
bit coordinate is either kept or flipped; different coordinates
cannot be exchanged because they have different colors. The
valuation vertices of a site are therefore precisely all
valuations in these labels. For a square, translating the
homogeneous equation by its coordinate flips changes only the
right-hand side, to some bit $\beta_B$. Its valuation vertices
are exactly the legal valuations of
Equation~\eqref{eq:square} with that bit.

No agreement between the flips used in different local
isomorphisms is being assumed. What is shared is the single
labeling already assigned to each target bit pair, so every
block uses the same coordinates for its incidences. Combining
the resulting block descriptions, the exclusion of other base
edges, and the private leaves gives $H\cong G(\beta)$.
\end{proof}

It remains to distinguish those shifts that merely relabel the
graph. A \emph{global offset assignment} gives a vector in
$\F^{d_t}$ to every horizontal name $x_{i,t,p,p'}$ and a scalar
in $\F$ to every vertical name $y_{i,\ell,p,q}$, with the same
value used in all blocks containing that name. Such an assignment
solves the shifted constraints if it satisfies
Equation~\eqref{eq:square} for every square; sites impose no
additional condition.

\begin{proposition}[Isomorphisms and global offsets]
\label{prop:isomorphism}
For every choice of square bits $\beta$, the graphs $G(0)$ and
$G(\beta)$ are isomorphic if and only if there is a global offset
assignment solving all shifted square constraints.
\end{proposition}

\begin{proof}
An isomorphism preserves degrees, hence every named base-color
class. Its action on each bit pair is a fixed flip. Assemble
these flips into an offset vector for each horizontal name and
an offset scalar for each vertical name. Each block of $G(0)$
has its all-zero valuation vertex. Its image is adjacent to
the bit-value vertices specified by these same offsets, so
its valuation is exactly their restriction to that block.
Legality of the image yields Equation~\eqref{eq:square} with
right-hand side $\beta_B$. Because a bit pair is shared by all
blocks using it, the resulting offsets form one global assignment.

Conversely, suppose such an assignment exists. Translate each
bit pair by its offset, and translate each block valuation by
the offsets on that block's coordinates. Linearity and
Equation~\eqref{eq:square} make the latter a bijection from the
homogeneous legal valuation set to the shifted legal set.
Shared offsets preserve all valuation--bit incidences and
nonincidences. The resulting color-preserving base-graph
isomorphism extends to the private leaves, whose numbers
depend only on the preserved base color.
\end{proof}

Proposition~\ref{prop:mate} reduces the analysis of every
equivalent graph to a list of shifted equations, while
Proposition~\ref{prop:isomorphism} identifies precisely which
lists admit an isomorphism. The next step is to extract
constraints on these shifts from the bijective game and to
express global solvability in terms of full addresses.
